\documentclass[10pt,a4paper]{amsart}
\usepackage[margin=19mm]{geometry}
\usepackage{amsmath,amssymb,mathtools}
\usepackage[scr=rsfs,cal=euler]{mathalfa}
\usepackage{xfrac}
\usepackage[unicode,hidelinks,bookmarks=false]{hyperref}
\newcommand{\ie}{i.e.\ }
\newcommand{\cf}{cf.\ }
\newcommand{\resp}{respectively\ }
\newcommand{\Subsection}{\subsection}
\providecommand{\nobreakdash}{\nobreak\hbox{-}}
\setlength{\emergencystretch}{2em}
\multlinegap0pt
\usepackage{bm}
\usepackage{bbm}
\usepackage{dsfont}
\usepackage{xspace}
\usepackage{cancel}
\usepackage{mathtools}
\usepackage{amsthm}
\makeatletter
\@ifundefined{theorem}{\newtheorem{theorem}{Theorem}}{}
\@ifundefined{lemma}{\newtheorem{lemma}[theorem]{Lemma}}{}
\makeatother
\def\Ch{\mathrm{Ch}}
\def\Ind{\mathrm{Ind}}
\def\del{\mathrm{del}}
\def\lk{\mathrm{lk}}
\title[On the Vertex Decomposability of $r$-Independence Complexes ]{On the Vertex Decomposability of $r$-Independence Complexes of Trees}
\author{Rutuja Sawant}
\date{Source: arXiv:2605.25150v1 (CC BY 4.0)}
\begin{document}
\maketitle

\noindent\emph{Source and licence.} On the Vertex Decomposability of $r$-Independence Complexes of Trees. Version arXiv:2605.25150v1 (CC BY 4.0), distributed under the Creative Commons Attribution 4.0 International licence, as recorded in the arXiv licence metadata for this exact version and shown on the abstract page https://arxiv.org/abs/2605.25150v1. Full rights notice: licenses/arxiv-2605.25150-CC-BY-4.0.txt.\par\smallskip

\noindent\textit{Editorial scope.} Counted unit: the Lemma whose source label is \texttt{lem: existence of shedding vertex in link} in the article (source0.tex lines 621--623, source label \texttt{lem: existence of shedding vertex in link}), delivered together with the article's own complete original proof of it (source0.tex lines 624--847, one \texttt{proof} environment of 224 lines). No mathematics is written, repaired or re-derived: statement and proof are the article's own text line for line. Required context retained uncounted: lem: all shedding vertices (lines 445--450). Every definition, display and label of those lines is retained unchanged. Omitted from this package and not redistributed: 1--444, 451--620, 848--1056. Nothing inside a retained excerpt is omitted or altered: every retained body is delivered exactly as the article's lines are written, so no omission receipt of any kind accompanies this package. Citations: the retained excerpts carry no citation command at all, so no bibliography entry of the article is transcribed and no reference is redistributed. Reference adaptation: none was needed and none was made; every reference inside the delivered unit resolves inside it, so no unresolved reference or citation is produced. Printed slips kept literal: no printed word, symbol, hypothesis or number of the counted statement or of its proof was altered. Layout and class: the article's own class is \texttt{amsart} with packages including fontenc, microtype, textcomp, fbb, geometry, amsmath, amssymb, amsthm, amscd, gensymb, graphicx, etoolbox, booktabs, stackrel; the wrapper is the lane's pinned \texttt{amsart} editorial wrapper at 10pt on A4, which restates the theorem-like environments the retained text needs and adds the article's own macro definitions that the retained text needs (\texttt{\textbackslash Ch}, \texttt{\textbackslash Ind}, \texttt{\textbackslash del}, \texttt{\textbackslash lk}), with extra wrapper packages (bm, bbm, dsfont, xspace, cancel, mathtools). The delivered numbering is the unit's own. Assets: no retained span contains a figure environment, a graphics-inclusion command, a PDF-page inclusion, a table or a TikZ picture; no image file is copied into this package, the package declares no asset dependency and no image file is redistributed with it. Licence: version arXiv:2605.25150v1 of this article is distributed under the Creative Commons Attribution 4.0 International licence, as recorded in the arXiv licence metadata for this exact version and shown on the abstract page https://arxiv.org/abs/2605.25150v1. A full rights notice is delivered at licenses/arxiv-2605.25150-CC-BY-4.0.txt. Characters: every retained excerpt is pure ASCII, so the delivered engine, XeLaTeX with its default Latin Modern text font, reports no missing character.
\begin{lemma}\label{lem: all shedding vertices}
Let $T$ be a tree with $|V(T)| \ge r+1$ and $r \ge 1$. A vertex $v \in V(T)$ is a shedding vertex of $\Ind_r(T)$ if and only if, when $T$ is viewed as rooted at $v$, there exists a subset $A \subseteq \mathrm{Ch}(v)$ such that $|V(T_u)| \le r$ for every $u \in A$ and 
\[
\Big|\Big(\bigsqcup_{u \in A} V(T_u)\Big) \cup \{v\}\Big| \ge r+1.
\]
\end{lemma}

\begin{lemma}\label{lem: existence of shedding vertex in link}
Let $T$ be a tree with $|V(T)| \ge r+1$ and $r \ge 1$. Let $v$ be a shedding vertex of $\Ind_r(T)$. Then, for every connected subgraph $C$ of $T$ containing $v$ and satisfying $|C| \le r-1$, the complex $\lk_{\Ind_r(T)}(V(C))$ has a shedding vertex $w$ such that the graph $C\cup\{w\}$ is connected.
\end{lemma}

\begin{proof}
Let $\Delta = \lk_{\Ind_r(T)}(C)$, and root the tree $T$ at $v$. Throughout, we do not distinguish between a vertex set and the subgraph it induces; in particular, we identify a subgraph with its vertex set when no ambiguity arises.

Since $v$ is a shedding vertex of $\Ind_r(T)$, Lemma~\ref{lem: all shedding vertices} yields a subset $A\subseteq \Ch(v)$ such that $|V(T_u)|\le r$ for every $u\in A$ and
\[
\left|\left(\bigsqcup_{u\in A}V(T_u)\right)\cup\{v\}\right|
\ge r+1.
\]
Set
\[
S=\bigsqcup_{u\in A}V(T_u).
\]
Then $|S|\ge r$ and the induced subgraph $T[S \cup \{v\}]$ is connected of size at least $r+1$. Since $|C|\le r-1$, the set $S\setminus C$ is non-empty. Choose
\[
w\in S\setminus C
\]
so that the rooted subtree \(T_w\) has the largest vertex set among all rooted subtrees \(T_x\) with \(x \in S \setminus C\).

We first show that $\operatorname{parent}(w)\in C$. Suppose $\operatorname{parent}(w) \notin C$. Then
\[
\operatorname{parent}(w)\in S\setminus C.
\]
Since $T_w$ is a proper rooted subtree of $T_{\operatorname{parent}(w)}$, we obtain
\[
|V(T_{\operatorname{parent}(w)})|
>
|V(T_w)|,
\]
contradicting the maximality of $T_w$. Hence $\operatorname{parent}(w)\in C$.

As $C$ is connected and $\operatorname{parent}(w)\in C$, it follows that the subgraph $C\cup\{w\}$ is connected.

We claim that \(w\) is a shedding vertex of \(\Delta\). By definition, it suffices to show that no facet of \(\del_\Delta(w)\) is a face of \(\lk_\Delta(w)\).

Let \(F\) be a facet of \(\del_\Delta(w)\). We will show that that $F\cup\{w\}\notin\Delta$.

We first show
\[
V(T_w)\setminus\{w\}\subseteq F.
\]

Let \(y\) be a descendant of \(w\), i.e., \(y \in V(T_w)\setminus\{w\}\). Suppose, for contradiction, that \(y\notin F\).
We show that
\[
F\cup\{y\}\in\del_\Delta(w),
\]
contradicting the fact that $F$ is a facet of $\del_{\Delta}(w)$.

Since \(y\) is a descendant of \(w\), every unique path from \(y\) to a vertex outside \(T_w\) passes through \(w\). Because $w\notin F\cup C$, it follows that in the induced subgraph
\[
T[F\cup C\cup\{y\}],
\]
the vertex \(y\) lies in a connected component entirely contained in \(T_w \setminus w\).

Moreover,
\[
|V(T_w)|\le r,
\]
because \(w\in S\) and every rooted subtree corresponding to a vertex in \(A\) has at most
\(r\) vertices. Hence every connected component of
\[
T[F\cup C\cup\{y\}]
\]
has at most \(r\) vertices. Therefore
\[
F\cup C\cup\{y\}\in\Ind_r(T),
\]
which implies that $F\cup\{y\}$ is a face of $\Delta$. Since $w\notin F\cup\{y\}$,
we obtain
\[
F\cup\{y\}\in\del_\Delta(w).
\]
But this properly contains the facet \(F\), a contradiction.

Therefore every descendant of \(w\) belongs to \(F\). Equivalently,
\[
V(T_w)\setminus\{w\}\subseteq F.
\]

We now consider two cases.

\medskip
\noindent
\textbf{Case 1.}
Assume that
\[
S\setminus(C\cup T_w)\subseteq F.
\]

Together with
\[
V(T_w)\setminus\{w\}\subseteq F,
\]
this implies that every vertex of \(S\setminus\{w\}\) belongs to \(F\cup C\). Recall, a shedding vertex $v \in C$. Hence
\[
S\cup\{v\}\subseteq F\cup C\cup\{w\}.
\]

Therefore the induced subgraph $T[F\cup C\cup\{w\}]$ contains the connected subgraph $T[S\cup\{v\}]$, which has at least \(r+1\) vertices. Consequently, $F\cup C\cup\{w\}$ is not a face of $\Ind_r(T)$, and hence
\[
F\cup\{w\}\notin\Delta.
\]

\medskip
\noindent
\textbf{Case 2.}
Assume now that there exists
\[
x\in S\setminus(C\cup T_w)
\]
such that $x\notin F$. Since \(F\) is a facet of \(\del_\Delta(w)\), the set $F\cup\{x\}$
cannot be a face of \(\del_\Delta(w)\); otherwise it would properly contain the facet \(F\). Hence
\[
F\cup\{x\}\notin\Delta.
\]
Because
\[
\Delta=\lk_{\Ind_r(T)}(C),
\]
it follows that the set $F\cup C\cup\{x\}$ is not a face of $\Ind_r(T)$.

Therefore the induced subgraph
\[
T[F\cup C\cup\{x\}]
\]
has a connected component \(C'\) satisfying
\[
|V(C')|\ge r+1.
\]

As $F$ is a facet of $\del_{\Delta}(w)$, it follows that
\[
F\cup C\in\Ind_r(T),
\]
every connected component of $T[F\cup C]$ has size at most \(r\). Thus the large connected component \(C'\) appears only after adjoining
the vertex \(x\), and therefore
\[
x\in V(C').
\]

Now $x\in S$,
so there exists \(u\in A\) such that $x\in V(T_u)$.
Since
\[
|V(T_u)|\le r
\]
while
\[
|V(C')|\ge r+1,
\]
the component \(C'\) cannot be a subgraph of \(T_u\). Hence \(C'\) contains a
vertex outside \(T_u\).

Because \(C'\) is connected, the unique path from \(x\) to a vertex outside \(T_u\) lies inside
\(C'\). This path must pass through \(u\), and therefore through the root \(v\). Hence
\[
v\in V(C').
\]

Since \(v \in C\) and \(v \in C'\), and \(C'\) is a connected component of \(T[F \cup C \cup \{x\}]\), the connected subgraph $C$ lies entirely in \(C'\). Hence $C$ is a subgraph of $C'$.

Now consider the graph obtained from \(C'\) by deleting the vertices of $T_x$.
Since every vertex of \(T_x \setminus \{x\}\) is connected to the rest of \(T\) only through \(x\), removing \(V(T_x)\) does not disconnect any path in \(C' \setminus V(T_x)\). Hence
\[
C'' := C' \setminus V(T_x)
\]
is connected.

We claim that \(C \subseteq C''\). Since \(C \subseteq C'\), it suffices to show that \(C \cap V(T_x) = \emptyset\).

Suppose, for contradiction, that there exists a vertex \(y \in C \cap V(T_x)\). Since \(T_x \setminus \{x\}\) is attached to the rest of the tree only through \(x\), any path in \(T\) from \(y\) to a vertex outside \(T_x\) must pass through \(x\). Because \(C\) is connected and contains such a vertex outside \(T_x\) (as \(x \notin C\)), this would force \(x \in C\), contradicting \(x \in S \setminus C\). Hence \(C \cap V(T_x) = \emptyset\).

Therefore no vertex of \(C\) is removed when passing from \(C'\) to \(C'' = C' \setminus V(T_x)\), and thus
\[
C \subseteq C''.
\]

Next we attach the rooted subtree \(T_w\) to $C''$ through the $\operatorname{parent}(w)$. Since $\operatorname{parent}(w)\in C$ which is a subgraph of  $C''$,
the root \(w\) is adjacent to a vertex of \(C''\). Therefore
\[
C''\cup T_w
\]
is connected.

Moreover, by the maximal choice of \(w\),
\[
|V(T_w)|\ge |V(T_x)|.
\]
Hence
\[
|V(C'') \cup V(T_w)|
\ge
|V(C'') \cup V(T_x)| \ge r+1.
\]

Finally, since
\[
V(T_w)\setminus\{w\}\subseteq F,
\]
every vertex of \(T_w\) belongs to \(F\cup\{w\}\). Consequently,
\[
C''\cup T_w
\]
is a connected subgraph of
\[
T[F\cup C\cup\{w\}].
\]

Therefore
\[
T[F\cup C\cup\{w\}]
\]
contains a connected subgraph with at least \(r+1\) vertices. Hence
\[
F\cup C\cup\{w\}\notin\Ind_r(T),
\]
and thus
\[
F\cup\{w\}\notin\Delta.
\]

Thus no facet of \(\del_\Delta(w)\) is a facet of \(\lk_\Delta(w)\). Therefore \(w\) is a
shedding vertex of \(\Delta\).
\end{proof}

\end{document}
